\section{Evaluation}
\subsection{Experimental Setup}
% 我们的实验一共在两个Cluster上进行。在Cluster A上，我们在72F CPU上评估了MFold的性能。同时，我们将Cluster B上运行的FastFold{\color{red}~\cite{}}作为我们的基线。每个Cluster的具体配置如表~\ref{tab:cluster}所示。
We evaluate our optimizations on two distinct computing environments. The primary environment is a CPU cluster that serves as the optimization target of this work. The cluster is equipped with two 72F processors, each consisting of 8 NUMA nodes. Unless otherwise stated, all reported CPU results are obtained on a single 72F processor, which allows us to highlight the effectiveness of our optimizations without the influence of inter-socket communication overhead. This platform provides not only high computational throughput but also specialized matrix computation units, which we exploit to accelerate AlphaFold2 inference.

For comparison, we benchmark against FastFold, a state-of-the-art GPU-accelerated implementation of AlphaFold2. FastFold experiments are conducted on a server equipped with two Intel(R) Xeon(R) Gold 6336Y processors and a single NVIDIA A100 GPU. This configuration represents the mainstream deployment setting for AlphaFold2 inference on GPUs and provides a strong baseline to evaluate the competitiveness of our CPU-based optimizations.

Our evaluation dataset consists primarily of protein sequences from CASP14, which is widely adopted as a standard benchmark for protein structure prediction. Using CASP14 ensures that our results are comparable with prior studies and faithfully reflect AlphaFold2's predictive capability under realistic inference workloads.

It is worth noting that our optimized implementation achieves inference performance on a single 72F processor that matches or even surpasses that of an A100 GPU running FastFold, highlighting the potential of CPU-based acceleration with matrix units and quantization techniques.
\begin{table}
% Table generated by Excel2LaTeX from sheet 'Sheet1'
\begin{tabular}{|c|c|c|}
    \hline
    Cluster & A   & B \bigstrut\\
    \hline
    \hline
    CPU & 72F & Intel(R) Xeon(R) Gold 6336Y \bigstrut\\
    \hline
    Core & more than 256 & 24 \bigstrut\\
    \hline
    Memory &     & 512GB DDR4 \bigstrut\\
    \hline
    GPU & -   & A100 \bigstrut\\
    \hline
    System &     & Ubuntu 22.04 Linux 5.15.0 \bigstrut\\
    \hline
    \end{tabular}%
    
\caption{Cluster Configuration}
\label{tab:cluster}%
\end{table}

\subsection{End to End Performance Comparison}
% 我们从CASP14中选择了从长度32到2180的不同的数据进行了测试。对比结果如图~\ref{fig:endtoend}所示。

\subsection{Attention Kernel Performance}
% 由于我们的MFold实现与FastFold并未共享同一pipeline，为了更加充分的展示加速的效果，我们对比了MFold的fp16、int8精度以及Fastfold的Attention模块的性能。我们使用了从32~4096的不同长度的序列进行测试。对比结果如图~\ref{fig:attention}所示。